\chapter{Evaluation}
\label{ch:evaluation}

Earlier chapters introduced the release's quantitative evidence
piecemeal, in service of whatever argument was at hand: Chapter~\ref{ch:procint}
cited the declaration catalog's proven/stated split while presenting
\textsf{procint}'s module architecture; Chapter~\ref{ch:correspondence}
cited the correspondence rail's status while presenting the D1 obligation;
Chapter~\ref{ch:empirical} cited rslab's receipted test and benchmark
counts while keeping them apart from kernel-checked standing. This
chapter consolidates that evidence into a single evaluation treatment,
following the structure of \texttt{paper/main.tex}'s own Evaluation
section: it explains what each rendered table measures and how the
number was computed, then defers to the rendered fragment itself for the
number, rather than restating any count, hash, or status literal by hand
in this chapter's prose. As Chapter~\ref{ch:mfact} argued and
Chapter~\ref{ch:procint} exercised, a number typed into prose is exactly
the kind of unreceipted assertion this thesis's manufacturing discipline
exists to avoid; every quantitative claim below is either a macro from
\texttt{release\_macros.tex} or an \verb|\input| of a ledgered fragment
rendered from \texttt{release-manifest.json} or \texttt{quadrature.json}.

\section{Formal Standing Evaluation}
\label{sec:eval-formal}

The first table this chapter reports is rendered directly from the
release manifest of the certified run, \texttt{release-manifest.json},
and from nothing else: no cell is written by hand, and the generation
pipeline fails closed, so that if the manifest is absent or stale the
fragment does not render and the surrounding document does not build for
release. What the table measures is the release's own accounting of its
declaration catalog: how many of the \TotalDecls\ declarations recorded
at release time reached kernel-admitted, axiom-audited proven status
(\KernelTheorems), how many remain formally stated with an open proof
obligation (\StatedCount), and the release-wide \texttt{sorry} and
unauthorized-\texttt{admit} counts (\SorryCount\ and \AdmitCount\
respectively, both required to be zero for \CertifiedStatus\ to read
PASS). These are not independently re-derived by this thesis; they are
the same manifest values \texttt{release\_macros.tex} exposes as
\LaTeX\ macros and Chapter~\ref{ch:procint} has already cited in context.

\input{../paper/evaluation}

\paragraph{Final status.} Table~\ref{tab:mfact-final-status} in
Chapter~\ref{ch:mfact}, generated from the release manifest and the
standing records that accompany it, aggregates the individual gate
outcomes --- axiom audit, negative fixtures, certification, and standing
quadrature --- into the single-row summary that \texttt{just status}
(Chapter~\ref{ch:ai-development}) reads back at the command line; it is
reproduced there rather than repeated here.

\section{Standing Quadrature}
\label{sec:eval-quadrature}

The release is not evaluated only by whether the Lean packages build.
\texttt{paper/main.tex} is explicit that it is evaluated by whether the
declaration catalog, the admitted Lean corpus, the process evidence, the
release manifest, and the paper's own claims form a \emph{closed standing
graph}: the Standing Quadrature check rejects orphan declarations,
unsupported claims, untraced artifacts, and any divergence between the
manifest and the paper. This closure property is itself kernel-checked
rather than asserted by a separate script running outside Lean's trust
boundary: a generated witness module, \texttt{ProcInt.Release.Quadrature},
carries the name-sorted proven surfaces of the catalog, the manifest, and
the axiom audit as literal lists and proves their pairwise equality by
\texttt{rfl} --- so that any orphan declaration on any one side makes the
corresponding lists differ syntactically and the kernel refuses the
witness outright, rather than the check silently passing on a mismatch.
The same witness module encodes the manufacturing run itself as a
\textsf{procint} process trace and proves, using the corpus's own
\texttt{DeclareConstraint} semantics from Chapter~\ref{ch:procint}, that
the run conforms to its declared process model: \textsf{procint} is, in
this one specific and kernel-checked sense, used to model the process by
which \textsf{mfact} manufactured \textsf{procint}. The table below is
rendered from \texttt{quadrature.json}; three negative controls --- a
corrupted evaluation number, a claim with no evidence, and a catalog
declaration with no rendered symbol --- are each engineered to cause a
typed refusal rather than a silent pass, in the same spirit as the
negative fixtures Chapter~\ref{ch:procint} and Chapter~\ref{ch:mfact}
describe for the Lean corpus itself.

% RENDERED by ggen sync from the quadrature graph (quadrature-pack).
% Do not edit: regenerate via `just standing-quadrature`.
\begin{table}[h]
\centering
\small
\begin{tabular}{lrl}
\toprule
Surface pair & Count & Result \\
\midrule
TTL to Lean & 160 & PASS \\
Lean to Audit & 160 & PASS \\
Audit to Manifest & 160 & PASS \\
Manifest to Paper & 6 & PASS \\
Artifact to Process Event & 215 & PASS \\
Claim to Evidence & 5 & PASS \\
\bottomrule
\end{tabular}
\caption{Standing Quadrature: PASS over
160 proven declarations, 5 traced paper
claims, and 6 evaluation numbers. The closure is
kernel-checked by the rendered witness
\texttt{ProcInt.Release.Quadrature}.}
\label{tab:quadrature}
\end{table}


\section{Correspondence Status}
\label{sec:eval-correspondence}

Chapter~\ref{ch:correspondence} developed the Aeneas/Charon correspondence
rail in detail, including the D1 token-replay-counts obligation, the
guardrail requiring a correspondence theorem to quantify over its actual
extracted declaration rather than a same-shaped hand-written stand-in, and
that rail's own generated evaluation table (rendered by the correspondence
builder from its receipt --- a charon/aeneas pipeline receipt, a
\texttt{lake build} of the extracted module against \textsf{procint}, and
an axiom-print check on the resulting theorem --- never written by hand),
reproduced there rather than repeated here.

\section{Falsifier and Valid Objection Surface}
\label{sec:eval-falsifier}

The evaluation above answers ``what does the release currently claim.''
\texttt{paper/main.tex}'s Falsifier and Valid Objection Surface section
answers a complementary question: what would it take for one of those
claims to be wrong, and is that possibility engineered into the release
gate rather than left as an informal hope. A valid objection to the
release, in the sense that section defines, must exhibit at least one of
eight closed failure classes: a declared admitted theorem contains
\texttt{sorry}; a theorem depends on an axiom outside the authorized set;
a manifest hash fails to match its artifact; a negative fixture passes
when it is designed to fail; a generated evaluation number fails to match
the manifest that is supposed to have produced it; a formally stated
declaration is described, anywhere in the paper or its fragments, as
proven; the catalog-to-Lean rendering step changes the theorem statement
it was supposed to project; or the release cannot be replayed under its
pinned dependencies.

These eight classes are deliberately enumerated as \emph{closed} failure
classes rather than left as an open-ended attack surface: each one is
engineered into a specific gate, a typed refusal, or a non-admitted
state, so that the release gate's design goal is that any occurrence of
one of these cases prevents admission or produces a typed refusal, rather
than passing silently through to a certified release. After certification,
no constructor of this valid-objection surface is inhabited within the
declared boundary --- the surface is empty in the sense that its
negative controls, described alongside Section~\ref{sec:eval-quadrature}'s
quadrature controls and Chapter~\ref{ch:procint}'s negative fixtures,
each demonstrate the corresponding gate actually refusing on a poisoned
copy, rather than the emptiness being asserted without a check that could
have failed. It is worth restating the scope of this claim precisely, in
the same register \texttt{paper/main.tex} insists on: a local objection
to one proof, one citation, or one template is not automatically a
refutation of the release's mathematical-manufacturing claim as a whole.
It becomes one only if it crosses the admission boundary and invalidates
a specific, claimed receipt --- and the eight classes above are exactly
the enumerated ways of doing that, not a suggestive sample of a larger,
unstated set. This same enumeration is the basis for
Chapter~\ref{ch:discussion}'s discussion of what admission does, and does
not, check.
